%! TEX root = ../note-dqm-transform.tex

\section{Conditional Expectations}
% \label{sec--conditioning}

\subsection{Auxiliary Lemmas and their Proofs}

\begin{lemma}
\label{lem--nonneg-f-equal-cond-exp-pos-prodf}
Let \(f\) be a non-negative integrable function on the probability space
\((\mathbf{X}, \mathcal{B}, \mu)\).
Let \(\mathcal{B}_{0}\) be a sub-\(\sigma\)-algebra of \(\mathcal{B}\), and
let \(B = \left\{ \mu \left[ f \middle| \mathcal{B}_{0} \right] > 0 \right\}\).
Then \(f = f \cdot \mathbf{1}_{B}\)
\(\mu\)-a.e.
\end{lemma}

\begin{proof}[Proof of \Cref{lem--nonneg-f-equal-cond-exp-pos-prodf}]
Deduce that \(B \in \mathcal{B}_{0}\) and hence so is \(A := \mathbf{X}
\setminus B = \left\{ \mu \left[ f \middle| \mathcal{B}_{0}
\right] = 0 \right\}\), where the last equality follows from \(f \geq 0\)
\(\mu\)-a.e.
By definition of \(\mu \left[ f \middle| \mathcal{B}_{0} \right]\),
\begin{equation*}
  \mu \left[ f \mathbf{1}_{A} \right] = \mu \left[ \mu \left[ f \middle|
  \mathcal{B}_{0} \right] \mathbf{1}_{A} \right] = 0,
\end{equation*}
where the last equality follows by definition of \(A\).
Since \(f \mathbf{1}_{A} \geq 0\), we conclude that \(f \mathbf{1}_{A} = 0\)
\(\mu\)-a.e.
Therefore, \(f = f \mathbf{1}_{A} + f \mathbf{1}_{B} = f \mathbf{1}_{B}\)
\(\mu\)-a.e.
\end{proof}

\begin{lemma}
\label{lem--nonneg-f-cond-exp-null-subsets}
Let \((\mathbf{X}, \mathcal{B}, \mu)\) be a probability space,
\(\mathcal{B}_{0}\) be a sub-\(\sigma\)-algebra of \(\mathcal{B}\),
and \(f\) be an integrable function with \(f \geq 0\) \(\mu\)-a.e.
Then \(\mathbf{1} \left\{ \mu \left[ f \middle| \mathcal{B}_{0}
\right] = 0 \right\} \leq \mathbf{1} \{f = 0\}\) \(\mu\)-a.e.
\end{lemma}

\begin{proof}[Proof of \Cref{lem--nonneg-f-cond-exp-null-subsets}]
Let \(\psi = \mu \left[ f \middle| \mathcal{B}_{0}
\right]\) and note that \(\left\{ \psi = 0 \right\} \in \mathcal{B}_{0}\) and
\(\psi \cdot \mathbf{1} \{\psi = 0\} = 0\) everywhere.
Therefore,
\(0 = \mu \left[ \psi \cdot \mathbf{1} \left\{ \psi = 0 \right\} \right] = \mu
\left[ f \cdot \mathbf{1} \left\{ \psi = 0 \right\} \right]\).
Deduce that a \(\mathbf{1} \{f > 0\} \mathbf{1} \left\{ \psi = 0 \right\} = 0\)
almost everywhere \(\mu\).
Therefore, by \(\mu\) almost everywhere non-negativity of \(f\),
\(\mathbf{1} \left\{ \psi = 0 \right\} = \mathbf{1} \{f = 0\} \mathbf{1} \left\{
\psi = 0 \right\} \leq \mathbf{1} \{f = 0\}\) with the first equality holding
\(\mu\)-a.e.
\end{proof}

\begin{lemma}
\label{lem--sqrt-integral-inequality}
Let \(f_{1}, f_{2}, f_{3}\) be measurable functions on the probability space
\((\mathbf{X}, \mathcal{B}, \mu)\) such that \(f_{1}, f_{2} \geq 0\)
\(\mu\)-a.e, \(f_{1}, f_{2}, f_{3} \in \mathscr{L}_{1} (\mu)\) and
\(\frac{f_{1} - f_{2} - f_{3}}{\sqrt{f_{1}} + \sqrt{f_{2}}} \in \mathscr{L}_{2}
(\mu)\).
Let \(\mathcal{B}_{0}\) be a sub-\(\sigma\)-algebra of \(\mathcal{B}\), and
define
\begin{equation*}
  B = \left\{ \sqrt{\mu \left[ f_{1} \middle| \mathcal{B}_{0} \right]} +
  \sqrt{\mu \left[ f_{2} \middle| \mathcal{B}_{0} \right]} > 0 \right\}.
\end{equation*}
Then
\begin{equation*}
  \begin{split}
    & \left( \sqrt{\mu \left[ f_{1} \middle| \mathcal{B}_{0} \right]} -
    \sqrt{\mu \left[ f_{2} \middle| \mathcal{B}_{0} \right]} -
    \frac{\mu \left[ f_{3} \middle| \mathcal{B}_{0} \right]}{\sqrt{\mu \left[
    f_{1} \middle| \mathcal{B}_{0} \right]} + \sqrt{\mu \left[ f_{2} \middle|
    \mathcal{B}_{0} \right]}} \mathbf{1}_{B} \right)^{2} \\
    \leq
    & \, \mu \left[ \left( \sqrt{f_{1}} - \sqrt{f_{2}} -
    \frac{f_{3}}{\sqrt{f_{1}} + \sqrt{f_{2}}} \right)^{2} \middle|
    \mathcal{B}_{0} \right].
  \end{split}
\end{equation*}
\end{lemma}

\begin{proof}[Proof of \Cref{lem--sqrt-integral-inequality}]
Rewrite
\begin{equation*}
    \mu \left[ \left( \sqrt{f_{1}} - \sqrt{f_{2}} - \frac{f_{3}}{\sqrt{f_{1}} +
    \sqrt{f_{2}}} \right)^{2} \middle| \mathcal{B}_{0} \right] = \mu \left[
    \left( \frac{f_{1} - f_{2} - f_{3}}{\sqrt{f_{1}} + \sqrt{f_{2}}} \right)^{2}
    \middle| \mathcal{B}_{0} \right].
\end{equation*}
By \Cref{lem--integral-of-ratio-squared-inequality},
\begin{equation*}
  \left( \mu \left[ f_{1} - f_{2} - f_{3} \middle| \mathcal{B}_{0} \right]
  \right)^{2} \leq
  \mu \left[ \left( \frac{f_{1} - f_{2} - f_{3}}{\sqrt{f_{1}} + \sqrt{f_{2}}}
  \right)^{2} \middle| \mathcal{B}_{0} \right] \mu \left[ \left( \sqrt{f_{1}} +
  \sqrt{f_{2}} \right)^{2} \middle| \mathcal{B}_{0} \right].
\end{equation*}
Furthermore, by Cauchy-Schwarz,
\begin{equation*}
  \begin{split}
    \mu \left[ \left( \sqrt{f_{1}} + \sqrt{f_{2}} \right)^{2} \middle|
    \mathcal{B}_{0} \right] =
    & \, \mu \left[ f_{1} \middle| \mathcal{B}_{0} \right] + \mu \left[ f_{2}
    \middle| \mathcal{B}_{0} \right] + 2 \mu \left[ \sqrt{f_{1}
    f_{2}} \middle| \mathcal{B}_{0} \right] \\
    \leq
    & \, \mu \left[ f_{1} \middle| \mathcal{B}_{0} \right] + \mu \left[ f_{2}
    \middle| \mathcal{B}_{0} \right] + 2 \sqrt{\mu \left[ f_{1} \middle|
    \mathcal{B}_{0} \right] \mu \left[ f_{2} \middle| \mathcal{B}_{0} \right]}
    \\
    =
    & \, \left( \sqrt{\mu \left[ f_{1} \middle| \mathcal{B}_{0} \right]} +
    \sqrt{\mu \left[ f_{2} \middle| \mathcal{B}_{0} \right]} \right)^{2}.
  \end{split}
\end{equation*}
Combining the previous two displays,
\begin{equation*}
  \left( \mu \left[ f_{1} - f_{2} - f_{3} \middle| \mathcal{B}_{0} \right]
  \right)^{2} \leq
  \mu \left[ \left( \frac{f_{1} - f_{2} - f_{3}}{\sqrt{f_{1}} +
  \sqrt{f_{2}}} \right)^{2} \middle| \mathcal{B}_{0} \right]
  \left( \sqrt{\mu \left[ f_{1} \middle| \mathcal{B}_{0} \right]} + \sqrt{\mu
  \left[ f_{2} \middle| \mathcal{B}_{0} \right]} \right)^{2}.
\end{equation*}
Multiply both sides by \(\mathbf{1}_{B}\) to get
\begin{equation*}
  \left( \mu \left[ f_{1} - f_{2} - f_{3} \middle| \mathcal{B}_{0} \right]
  \right)^{2} \mathbf{1}_{B} \leq
  \mu \left[ \left( \frac{f_{1} - f_{2} - f_{3}}{\sqrt{f_{1}} +
  \sqrt{f_{2}}} \right)^{2} \middle| \mathcal{B}_{0} \right]
  \left( \sqrt{\mu \left[ f_{1} \middle| \mathcal{B}_{0} \right]} + \sqrt{\mu
  \left[ f_{2} \middle| \mathcal{B}_{0} \right]} \right)^{2} \mathbf{1}_{B}.
\end{equation*}
On \(B\), we can divide by the second quantity on the right hand side of the
inequality so that
\begin{equation*}
  \begin{split}
    \left( \frac{\mu \left[ f_{1} - f_{2} - f_{3} \middle| \mathcal{B}_{0}
    \right]}{\sqrt{\mu \left[ f_{1} \middle| \mathcal{B}_{0} \right]} +
    \sqrt{\mu \left[ f_{2} \middle| \mathcal{B}_{0} \right]}} \right)^{2}
    \mathbf{1}_{B} \leq
    & \, \mu \left[ \left( \frac{f_{1} - f_{2} - f_{3}}{\sqrt{f_{1}} +
    \sqrt{f_{2}}} \right)^{2} \middle| \mathcal{B}_{0} \right] \mathbf{1}_{B} \\
    \leq & \, \mu \left[ \left( \frac{f_{1} - f_{2} -
    f_{3}}{\sqrt{f_{1}} + \sqrt{f_{2}}} \right)^{2} \middle| \mathcal{B}_{0}
    \right].
  \end{split}
\end{equation*}
By \(f_{j} \geq 0\), \(\mu \left[ f_{j} \middle| \mathcal{B}_{0} \right] \geq
0\) for both \(j = 1, 2\).
Therefore,
\begin{equation*}
  \begin{split}
    \frac{\mu \left[ f_{1} - f_{2} - f_{3} \middle| \mathcal{B}_{0}
    \right]}{\sqrt{\mu \left[ f_{1} \middle| \mathcal{B}_{0} \right]} +
    \sqrt{\mu \left[ f_{2} \middle| \mathcal{B}_{0} \right]}}
    \mathbf{1}_{B} =
    & \, \left( \sqrt{\mu \left[ f_{1} \middle| \mathcal{B}_{0} \right]}
    - \sqrt{\mu \left[ f_{2} \middle| \mathcal{B}_{0} \right]} -
    \frac{\mu \left[ f_{3} \middle| \mathcal{B}_{0}
    \right]}{\sqrt{\mu \left[ f_{1} \middle| \mathcal{B}_{0} \right]} +
    \sqrt{\mu \left[ f_{2} \middle| \mathcal{B}_{0} \right]}} \right)
    \mathbf{1}_{B} \\
    =
    & \, \left( \sqrt{\mu \left[ f_{1} \middle| \mathcal{B}_{0} \right]}
    - \sqrt{\mu \left[ f_{2} \middle| \mathcal{B}_{0} \right]}  \right)
    \mathbf{1}_{B} - \frac{\mu \left[ f_{3} \middle| \mathcal{B}_{0}
    \right]}{\sqrt{\mu \left[ f_{1} \middle| \mathcal{B}_{0} \right]} +
    \sqrt{\mu \left[ f_{2} \middle| \mathcal{B}_{0} \right]}} \mathbf{1}_{B} \\
    =
    & \,
    \sqrt{\mu \left[ f_{1} \middle| \mathcal{B}_{0} \right]} - \sqrt{\mu \left[
    f_{2} \middle| \mathcal{B}_{0} \right]} - \frac{\mu \left[ f_{3} \middle|
    \mathcal{B}_{0} \right]}{\sqrt{\mu \left[ f_{1} \middle| \mathcal{B}_{0}
    \right]} + \sqrt{\mu \left[ f_{2} \middle| \mathcal{B}_{0} \right]}}
    \mathbf{1}_{B},
  \end{split}
\end{equation*}
where the last line follows from non-negativity of \(\mu \left[ f_{j} \middle|
\mathcal{B}_{0} \right]\) and the definition of \(B\).
Combining the above two displays, we get
\begin{equation*}
  \begin{split}
    & \left( \sqrt{\mu \left[ f_{1} \middle| \mathcal{B}_{0} \right]} -
    \sqrt{\mu \left[ f_{2} \middle| \mathcal{B}_{0} \right]} -
    \frac{\mu \left[ f_{3} \middle| \mathcal{B}_{0} \right]}{\sqrt{\mu \left[
    f_{1} \middle| \mathcal{B}_{0} \right]} + \sqrt{\mu \left[ f_{2} \middle|
    \mathcal{B}_{0} \right]}} \mathbf{1}_{B} \right)^{2} \\
    \leq & \, \mu \left[ \left( \frac{f_{1} - f_{2} -
    f_{3}}{\sqrt{f_{1}} + \sqrt{f_{2}}} \right)^{2} \middle| \mathcal{B}_{0}
    \right] \\
    \leq
    & \, \mu \left[ \left( \sqrt{f_{1}} - \sqrt{f_{2}} -
    \frac{f_{3}}{\sqrt{f_{1}} + \sqrt{f_{2}}} \right)^{2} \middle|
    \mathcal{B}_{0} \right].
  \end{split}
\end{equation*}
\end{proof}

\begin{lemma}
\label{lem--ratio-integrable-implication}
Let \(r \in [1, \infty)\) \(f, g\) be measurable functions on a probability
space \((\mathbf{X}, \mathcal{B}, \mu)\) with \((f / g) \in
\mathscr{L}_{r} (\mu)\).
Then
\begin{equation}
  \mu \{|f| > 0, g = 0\} = 0, \quad \text{and} \quad f \mathbf{1} \{g = 0\} = 0
  \quad \mu\text{-a.e}.
  \label{eqn--ratio-integrable-implication}
\end{equation}
\end{lemma}

\begin{proof}[Proof of \Cref{lem--ratio-integrable-implication}]
The second part of \eqref{eqn--ratio-integrable-implication} is an immediate
implication of the first part.
It suffices to prove the first part for the case \(r = 1\), since the same proof
applies with \(|f| / |g|\) replaced by \(|f|^{r} / |g|^{r}\).
We prove the first part of \eqref{eqn--ratio-integrable-implication} by
contrapositive.

Suppose that \(\mu \{|f| > 0, g = 0\} > 0\).
Then there is some \(\varepsilon > 0\) such that \(\mu \{|f| >
\varepsilon, g = 0\} > 0\), since otherwise \(\mu \{|f| > (1 / n), g
= 0\} = 0\) for every \(n \in \mathbb{N}\) and \(\mu \{|f| > 0, g = 0\} =
\lim_{n \to \infty} \mu \{|f| > (1 / n), g = 0\} = 0\), a contradiction.
But then \(f / g \not\in \mathscr{L}_{1} (\mu)\) because
\begin{equation*}
  \mu [|f| / |g|] \geq \mu [(|f| / |g|) \mathbf{1} \{|f| > \varepsilon, g = 0\}]
  \geq \varepsilon \mu [(1 / |g|) \mathbf{1} \{|f| > \varepsilon, g = 0\}] =
  \infty.
\end{equation*}
By contrapositive, \(\mu [|f| / |g|] < \infty\) implies \(\mu \{|f| > 0, g = 0\}
= 0\).
\end{proof}

\begin{lemma}
\label{lem--integral-of-ratio-squared-inequality}
Let \(f, g\) be measurable functions on a probability space \((\mathbf{X},
\mathcal{B}, \mu)\) with \(f \in \mathscr{L}_{1} (\mu)\) and \(g, (f / g) \in
\mathscr{L}_{2} (\mu)\).
Given any sub-\(\sigma\)-algebra \(\mathcal{B}_{0} \subseteq \mathcal{B}\),
\begin{equation}
  \left( \mu \left[ f \middle| \mathcal{B}_{0} \right] \right)^{2} \leq
  \mu \left[ \left( \frac{f}{g} \right)^{2} \middle| \mathcal{B}_{0} \right]
  \mu \left[ g^{2} \middle| \mathcal{B}_{0} \right].
  \label{eqn--integral-of-ratio-squared-inequality}
\end{equation}
\end{lemma}

\begin{proof}[Proof of \Cref{lem--integral-of-ratio-squared-inequality}]
By \Cref{lem--ratio-integrable-implication}, \((f / g) \in \mathscr{L}_{2}
(\mu)\) implies that \(f = f \mathbf{1} \{g \neq 0\}\) and
\((f / g) = (f / g) \mathbf{1} \{g \neq 0\}\).
To see this, note that it suffices to prove the first equality and so,
\begin{equation*}
  \begin{gathered}
    f = f \mathbf{1} \{g \neq 0\} + \underset{= 0 \
    \mu\text{-a.e.}}{\underbrace{f \mathbf{1} \{g = 0\}}} = f \mathbf{1} \{g
    \neq 0\} \quad \mu\text{-a.e}, \\
    \text{so that }
    \mu \left[ f \middle| \mathcal{B}_{0} \right] = \mu \left[ f
    \mathbf{1} \{g \neq 0\} \middle| \mathcal{B}_{0} \right] \text{ and }
    \mu \left[ (f / g)^{2} \middle| \mathcal{B}_{0} \right] = \mu \left[ (f
    / g)^{2} \mathbf{1} \{g \neq 0\} \middle| \mathcal{B}_{0} \right] \quad
    \mu\text{-a.e.}
  \end{gathered}
\end{equation*}
Then, \eqref{eqn--integral-of-ratio-squared-inequality} follows from the
Cauchy-Schwarz inequality:
\begin{equation*}
  \begin{split}
    (\mu \left[ f \middle| \mathcal{B}_{0} \right])^{2}
    =
    & \, (\mu \left[ f \mathbf{1} \{g \neq 0\} \middle| \mathcal{B}_{0}
    \right])^{2}
    = \left( \mu \left[ \frac{f}{g} g \cdot \mathbf{1} \{g \neq 0\} \middle|
    \mathcal{B}_{0} \right] \right)^{2} \\
    \leq & \, \mu \left[ \left( \frac{f}{g} \right)^{2} \cdot \mathbf{1} \{g
    \neq 0\} \middle| \mathcal{B}_{0} \right] \mu \left[ g^{2} \middle|
    \mathcal{B}_{0} \right].
  \end{split}
\end{equation*}
\end{proof}

%%% Local Variables:
%%% mode: LaTeX
%%% TeX-master: "../note-dqm-transform"
%%% End:
